\subsection{Time-dependent perturbation theory}
\begin{aside}{Why this subsection}
The product formula \eqref{eq:eint} needs commuting terms. When terms do not commute, or when boundary conditions are imposed, there is no closed formula, but $\Eint$ is still defined by \eqref{xeq:eint-def}. Perturbation theory expands the flow in powers of the interaction. From this expansion one reads
\begin{enumerate}[nosep,label=(\roman*)]
\item that $\Eint$ is well defined for all times (convergence),
\item approximations obtained by truncating the series order by order, and
\item the size of $\Eint$ for small $t$.
\end{enumerate}
\end{aside}

Split $A=A_1+B$, where $A_1$ is the principal term (for instance diffusion) and $B$, the remaining terms, is treated as a perturbation of $A_1$. The Dyson--Phillips series is
\begin{equation}\label{eq:dyson}
\begin{aligned}
S(t)&=\sum_{k\ge0}S^{(k)}(t),\\
S^{(0)}(t)&=e^{tA_1},\\
S^{(k+1)}(t)&=\int_0^te^{(t-s)A_1}\,B\,S^{(k)}(s)\dd s ,
\end{aligned}
\end{equation}
the same expansion that gives the interaction picture of quantum mechanics~\cite{dy}.

\begin{aside}{Reading the series}
\begin{itemize}[nosep]
\item $S^{(0)}=e^{tA_1}$ is the flow of $A_1$ alone;
\item $S^{(k)}$ collects the contributions in which $B$ acts $k$ times, at intermediate times $s\in(0,t)$, with the flow of $A_1$ in between.
\end{itemize}
The terms with $k\ge1$ are therefore the part of the flow in which $B$ acts on the evolution generated by $A_1$, that is, the interaction. (When $A_1$ and $B$ commute, the first-order term is
\[
\int_0^te^{(t-s)A_1}Be^{sA_1}\dd s=t\,Be^{tA_1},
\]
consistent with Section~\ref{sec:inside}.)
\end{aside}

\textbf{What it gives.} The series converges for all $t$ when $A_1$ generates an even-order dissipative semigroup (e.g.\ diffusion) and $B$ is a lower-order $L^\infty$ perturbation, also with Dirichlet conditions \led{P3-T1} \modulo{Miyadera--Voigt~\cite{en}}: so $\Eint$ is well defined and can be approximated by truncation even without a closed formula. Each $S^{(k)}$ is the order-$k$ part of the interaction. For small times, for boundary-induced interaction, with $k$ the order of contact of the datum with the boundary, $\|\Eint\|\sim t^{k/2+3/4}$ \led{P3-T3, OP-F5A} \proved\edmark{E10}: the smoother the datum is against the boundary (larger $k$), the faster the interaction becomes negligible as $t\to0$.

\subsection{Non-commuting terms}\label{sec:noncomm}
If $[A_1,A_2]\ne0$, then $e^{t(A_1+A_2)}\neq e^{tA_1}e^{tA_2}$; by Baker--Campbell--Hausdorff,
\[
e^{tA_1}e^{tA_2}=\exp\!\Big(t(A_1+A_2)+\tfrac{t^2}{2}[A_1,A_2]+O(t^3)\Big).
\]
\begin{proposition}[{(Lie-algebra dichotomy) \led{P1-T4, P1-T5} \proved[]}]
If the Lie algebra generated by the terms is solvable, an exact formula with finitely many scalar functions exists (Wei--Norman). Otherwise only a local expansion exists, and \textbf{no finite product of the flows of the individual terms} reproduces $e^{t(X+\varepsilon Y)}$.
\end{proposition}
\begin{takeaway}
For commuting terms the interaction is explicit and cheap to evaluate, but it cannot be made to vanish nor be computed stably in $H^m$; for non-commuting terms even the explicit form is lost in general. These are the linear-case limits against which the nonlinear extension of Section~\ref{sec:nonlin} must be read.
\end{takeaway}

%======================================================================
